\documentclass[a4paper,12pt,titlepage,final]{article}

\usepackage[T1,T2A]{fontenc}
\usepackage[utf8]{inputenc}
\usepackage[russian]{babel}

\usepackage{tikz}
\usepackage{pgfplots}
\usepgfplotslibrary{fillbetween}

\usepackage{geometry}
\usepackage{indentfirst}
\usepackage{amsmath}
\usepackage{amssymb}
\usepackage{array}
\usepackage{verbatim}

\geometry{a4paper,left=30mm,top=30mm,bottom=30mm,right=30mm}
\setcounter{secnumdepth}{0}
\pgfplotsset{compat=1.18}

\begin{document}

\begin{titlepage}
    \begin{center}
        {\small \sc Московский государственный университет\\
        имени М.~В.~Ломоносова\\
        Факультет вычислительной математики и кибернетики\\}

        \vfill

        {\Large \sc Отчет по заданию №6}\\
        ~\\
        {\large \bf <<Сборка многомодульных программ.\\
        Вычисление корней уравнений и определенных интегралов>>}\\
        ~\\
        {\large \bf Вариант 7}
    \end{center}

    \begin{flushright}
        \vfill
        {Выполнила:\\
        студентка 104 группы\\
        Конгаа~А.~Н.\\
        ~\\
        Преподаватель:\\
        Гуляев~Д.~А.}
    \end{flushright}

    \begin{center}
        \vfill
        {\small Москва\\2026}
    \end{center}
\end{titlepage}

\tableofcontents
\newpage

\section{Постановка задачи}

В данной работе требуется реализовать численные методы, позволяющие вычислить площадь плоской фигуры, ограниченной тремя кривыми:
\[
    y=f_1(x), \qquad y=f_2(x), \qquad y=f_3(x).
\]

Для варианта 7 заданы функции:
\[
    f_1(x)=\ln x,
\]
\[
    f_2(x)=-2x+14,
\]
\[
    f_3(x)=\frac{1}{2-x}+6.
\]

Для вычисления площади необходимо найти точки пересечения графиков функций, представить площадь фигуры как алгебраическую сумму определенных интегралов и вычислить эти интегралы с заданной точностью.

В работе используются следующие параметры.

\begin{center}
\begin{tabular}{|c|c|}
\hline
Параметр & Значение \\
\hline
Тип данных & \texttt{double} \\
\hline
Методы решения уравнений & деление пополам, комбинированный метод \\
\hline
Метод интегрирования & метод трапеций \\
\hline
Точности вычислений & $10^{-3}$, $10^{-5}$, $10^{-7}$ \\
\hline
Языки программирования & C, NASM32 \\
\hline
Соглашение вызова & \texttt{cdecl} \\
\hline
\end{tabular}
\end{center}

Функции $f_1$, $f_2$, $f_3$ реализованы на языке ассемблера NASM32 с использованием сопроцессора x87. Остальная часть программы реализована на языке C. Сборка выполняется с помощью \texttt{make}.

\newpage

\section{Математическое обоснование}

Для нахождения границ фигуры требуется найти три точки пересечения графиков. Для этого решаются уравнения:
\[
    f_1(x)=f_3(x), \qquad f_2(x)=f_3(x), \qquad f_1(x)=f_2(x).
\]

Функция $f_3(x)$ имеет разрыв в точке $x=2$, поэтому отрезки для поиска корней выбирались правее этой точки. На выбранных отрезках функции непрерывны, а разность соответствующих функций меняет знак, следовательно, метод деления пополам применим.

\begin{center}
\begin{tabular}{|c|c|c|}
\hline
Пересечение & Уравнение & Отрезок поиска \\
\hline
$x_{13}$ & $f_1(x)-f_3(x)=0$ & $[2.1;2.3]$ \\
\hline
$x_{23}$ & $f_2(x)-f_3(x)=0$ & $[4.0;4.5]$ \\
\hline
$x_{12}$ & $f_1(x)-f_2(x)=0$ & $[6.0;6.2]$ \\
\hline
\end{tabular}
\end{center}

На рисунке~\ref{plot1} изображены графики функций.

\begin{figure}[h]
\centering
\begin{tikzpicture}
\begin{axis}[
    axis lines=middle,
    xmin=1.8, xmax=6.5,
    ymin=0, ymax=11,
    width=13cm,
    height=8cm,
    samples=256,
    legend cell align=left
]
\addplot[green,thick,domain=1.9:6.5] {ln(x)};
\addlegendentry{$y=\ln x$}

\addplot[blue,thick,domain=1.9:6.5] {-2*x+14};
\addlegendentry{$y=-2x+14$}

\addplot[red,thick,domain=2.05:6.5] {1/(2-x)+6};
\addlegendentry{$y=\frac{1}{2-x}+6$}
\end{axis}
\end{tikzpicture}
\caption{Графики заданных функций}
\label{plot1}
\end{figure}

В результате численных расчетов были получены следующие абсциссы точек пересечения:
\[
    x_{13}\approx 2.191743425,\qquad
    x_{23}\approx 4.224744871,\qquad
    x_{12}\approx 6.096169674.
\]

На отрезке $[x_{13},x_{23}]$ верхней границей фигуры является график $f_3(x)$, а нижней границей является график $f_1(x)$. На отрезке $[x_{23},x_{12}]$ верхней границей фигуры является график $f_2(x)$, а нижней границей является график $f_1(x)$.

Следовательно, площадь фигуры вычисляется по формуле:
\[
S=
\int_{x_{13}}^{x_{23}}\left(f_3(x)-f_1(x)\right)\,dx+
\int_{x_{23}}^{x_{12}}\left(f_2(x)-f_1(x)\right)\,dx.
\]

Для поиска корней была выбрана точность $\varepsilon_1=10^{-5}$. Для вычисления интегралов была выбрана точность $\varepsilon_2=10^{-5}$. Эти значения меньше требуемой точности итоговой площади $\varepsilon=10^{-3}$ и дают устойчивый результат при вычислении площади.

\newpage

\section{Результаты экспериментов}

В результате проведенных вычислений были получены координаты точек пересечения графиков и площадь фигуры.

\begin{table}[h]
\centering
\begin{tabular}{|c|c|}
\hline
Величина & Значение \\
\hline
$x_{13}$ & 2.191743425 \\
\hline
$x_{23}$ & 4.224744871 \\
\hline
$x_{12}$ & 6.096169674 \\
\hline
$\displaystyle\int_{x_{13}}^{x_{23}}(f_3-f_1)\,dx$ & 9.1179 \\
\hline
$\displaystyle\int_{x_{23}}^{x_{12}}(f_2-f_1)\,dx$ & 2.1185 \\
\hline
Площадь фигуры $S$ & 11.23637645 \\
\hline
\end{tabular}
\caption{Результаты вычислений}
\label{table1}
\end{table}

Результат также был проверен при различных значениях точности.

\begin{table}[h]
\centering
\begin{tabular}{|c|c|}
\hline
$\varepsilon$ & Площадь фигуры \\
\hline
$10^{-3}$ & 11.236 \\
\hline
$10^{-5}$ & 11.23637 \\
\hline
$10^{-7}$ & 11.236376 \\
\hline
\end{tabular}
\caption{Зависимость результата от точности}
\end{table}

На рисунке~\ref{plot2} показана область, площадь которой вычисляется программой.

\begin{figure}[h]
\centering
\begin{tikzpicture}
\begin{axis}[
    axis lines=middle,
    xmin=1.8, xmax=6.5,
    ymin=0, ymax=11,
    width=13cm,
    height=8cm,
    samples=256,
    legend cell align=left,
    xticklabels={,,},
    yticklabels={,,}
]
\addplot[green,thick,domain=1.9:6.5,name path=A] {ln(x)};
\addlegendentry{$y=\ln x$}

\addplot[blue,thick,domain=1.9:6.5,name path=B] {-2*x+14};
\addlegendentry{$y=-2x+14$}

\addplot[red,thick,domain=2.05:6.5,name path=C] {1/(2-x)+6};
\addlegendentry{$y=\frac{1}{2-x}+6$}

\addplot[blue!20] fill between[of=C and A,soft clip={domain=2.1917:4.2247}];
\addplot[blue!20] fill between[of=B and A,soft clip={domain=4.2247:6.0962}];
\addlegendentry{$S\approx 11.2364$}

\addplot[dashed] coordinates {(2.1917,0) (2.1917,5.478)};
\addplot[color=black] coordinates {(2.1917,0)} node[label={-90:{\small 2.1917}}]{};

\addplot[dashed] coordinates {(4.2247,0) (4.2247,5.551)};
\addplot[color=black] coordinates {(4.2247,0)} node[label={-90:{\small 4.2247}}]{};

\addplot[dashed] coordinates {(6.0962,0) (6.0962,1.808)};
\addplot[color=black] coordinates {(6.0962,0)} node[label={-90:{\small 6.0962}}]{};
\end{axis}
\end{tikzpicture}
\caption{Плоская фигура, ограниченная графиками заданных уравнений}
\label{plot2}
\end{figure}

При уменьшении $\varepsilon$ значение площади стабилизируется. Это подтверждает корректность выбранных методов и выбранных отрезков поиска корней.

\newpage

\section{Структура программы и спецификация функций}

Программа состоит из нескольких модулей. Модули на языке C реализуют численные методы, обработку аргументов командной строки и вывод результатов. Модуль на языке NASM32 реализует заданные математические функции.

\begin{center}
\begin{tabular}{|c|p{10cm}|}
\hline
Файл & Назначение \\
\hline
\texttt{main.c} & основная программа, обработка ключей командной строки, вычисление площади \\
\hline
\texttt{root.c} & реализация методов решения уравнений \\
\hline
\texttt{integral.c} & реализация метода трапеций \\
\hline
\texttt{funcs.asm} & реализация функций $f_1$, $f_2$, $f_3$ на NASM32 \\
\hline
\texttt{funcs.h} & заголовочный файл с прототипами функций \\
\hline
\texttt{Makefile} & файл сборки программы \\
\hline
\end{tabular}
\end{center}

Разбиение программы на компоненты:
\[
\texttt{main.c}
\rightarrow
\begin{cases}
\texttt{root.c},\\
\texttt{integral.c},\\
\texttt{funcs.asm}.
\end{cases}
\]

Функция поиска корня имеет следующую спецификацию:
\begin{verbatim}
double root(
    double (*f)(double),
    double (*g)(double),
    double a,
    double b,
    double eps
);
\end{verbatim}

Функция ищет корень уравнения $f(x)-g(x)=0$ на отрезке $[a,b]$ с точностью $\varepsilon$.

Функция вычисления интеграла имеет следующую спецификацию:
\begin{verbatim}
double integral(
    double (*f)(double),
    double a,
    double b,
    double eps
);
\end{verbatim}

Функция вычисляет определенный интеграл функции $f$ на отрезке $[a,b]$ методом трапеций.

\newpage

\section{Сборка программы (Make-файл)}

Сборка программы выполняется с помощью \texttt{make}. По умолчанию собирается версия программы, использующая метод деления пополам. Для сборки версии с комбинированным методом используется макрос \texttt{COMBINED}.

\begin{verbatim}
make
make METHOD=COMBINED
make clean
\end{verbatim}

Зависимости между модулями программы:
\[
\texttt{main.c},\ \texttt{root.c},\ \texttt{integral.c},\ \texttt{funcs.asm}
\longrightarrow
\texttt{main}
\]

\newpage

\section{Отладка программы, тестирование функций}

Отладка программы проводилась поэтапно. Сначала проверялись ассемблерные функции $f_1$, $f_2$, $f_3$. После этого отдельно тестировались функции \texttt{root} и \texttt{integral}. Затем проверялась вся программа целиком.

Для проверки функции \texttt{root} использовалось уравнение с известным корнем:
\[
    x^2=4.
\]
На отрезке $[1;3]$ корень равен $x=2$.

Для проверки функции \texttt{integral} использовался интеграл:
\[
    \int_1^2 \ln x\,dx.
\]
Его точное значение:
\[
    2\ln2-1.
\]

В программе реализованы ключи командной строки:
\begin{itemize}
    \item \texttt{--help} --- вывод справки;
    \item \texttt{--points} --- вывод точек пересечения;
    \item \texttt{--iterations} --- вывод числа итераций;
    \item \texttt{--test} --- запуск тестов.
\end{itemize}

Полученные численные результаты совпали с аналитическими значениями с требуемой точностью.

\newpage

\section{Программа на Си и на Ассемблере}

Исходные тексты программы имеются в архиве, который приложен к этому отчету. Программа состоит из модулей на языке C и одного модуля на языке ассемблера NASM32. Взаимодействие между C и ассемблером выполняется по соглашению вызова \texttt{cdecl}.

Функции, реализованные на ассемблере:
\[
    f_1(x)=\ln x,
\]
\[
    f_2(x)=-2x+14,
\]
\[
    f_3(x)=\frac{1}{2-x}+6.
\]

\newpage

\section{Анализ допущенных ошибок}

Во время выполнения работы возникали ошибки, связанные с выбором отрезков поиска корней и с особенностями сборки многомодульной программы.

Основные ошибки:
\begin{itemize}
    \item неверный выбор отрезков поиска корней;
    \item попытка искать корень на отрезке, содержащем точку разрыва функции $f_3$;
    \item ошибки в именах ассемблерных функций при сборке под Windows и Linux;
    \item ошибки, связанные с соглашением вызова \texttt{cdecl};
    \item недостаточно малые значения $\varepsilon_1$ и $\varepsilon_2$ при вычислении площади.
\end{itemize}

После исправления ошибок программа стала корректно вычислять точки пересечения графиков и площадь плоской фигуры.

\newpage

\begin{raggedright}
\addcontentsline{toc}{section}{Список цитируемой литературы}
\begin{thebibliography}{99}
\bibitem{math} Ильин~В.\,А., Садовничий~В.\,А., Сендов~Бл.\,Х. Математический анализ. Т.\,1~---
Москва: Наука, 1985.
\end{thebibliography}
\end{raggedright}

\end{document}
